% gerado por scripts/gen_enem_tex.py — não editar
\voldef{1}{Números, Contagem, Estatística e Probabilidade}{Números, Contagem e Estatística}
\capdef{v1c01}{1}{1}{Conjuntos numéricos e operações}{Naturais, inteiros, racionais, reais; frações, decimais, potências, raízes, notação científica, ordem de grandeza, arredondamento, cálculo com números no cotidiano (contas, cardápios, troco, tempo).}
\capdef{v1c02}{1}{2}{Desigualdades, comparações e estimativas}{Comparar e ordenar números/frações/decimais, intervalos reais, inequações simples no contexto numérico, escolha da opção mais vantajosa por comparação, estimativas e razoabilidade de resultados.}
\capdef{v1c03}{1}{3}{Divisibilidade, fatoração, MDC e MMC}{Múltiplos e divisores, critérios de divisibilidade, números primos, fatoração, MDC, MMC, resto da divisão, ciclos e periodicidade, sistemas de numeração.}
\capdef{v1c04}{1}{4}{Razões e proporções}{Razão, proporção, densidade, velocidade média, consumo, taxas, razões de mistura, propriedades das proporções.}
\capdef{v1c05}{1}{5}{Grandezas proporcionais e regra de três}{Relação de dependência entre grandezas, diretamente e inversamente proporcionais, regra de três simples e composta, divisão proporcional.}
\capdef{v1c06}{1}{6}{Porcentagem}{Porcentagem, aumentos e descontos, variação percentual, aumentos/descontos sucessivos, pontos percentuais.}
\capdef{v1c07}{1}{7}{Juros e matemática financeira}{Juros simples, juros compostos, montante, financiamentos, parcelamentos, inflação, comparação de planos de pagamento.}
\capdef{v1c08}{1}{8}{Sequências e progressões}{Padrões numéricos e figurais, sequências, progressão aritmética, progressão geométrica, soma de termos.}
\capdef{v1c09}{1}{9}{Princípios de contagem}{Princípio fundamental da contagem, permutações, arranjos, combinações, anagramas, permutação com repetição, senhas e placas.}
\capdef{v1c10}{1}{10}{Representação e análise de dados}{Leitura e interpretação de tabelas, gráficos de barras, linhas, setores, pictogramas, histogramas, infográficos, inferências a partir de dados.}
\capdef{v1c11}{1}{11}{Medidas de tendência central}{Média aritmética, média ponderada, mediana, moda.}
\capdef{v1c12}{1}{12}{Medidas de dispersão}{Amplitude, desvio, desvio médio, variância, desvio padrão, regularidade/homogeneidade.}
\capdef{v1c13}{1}{13}{Probabilidade}{Probabilidade clássica, evento complementar, união, probabilidade condicional, eventos independentes, experimentos sucessivos.}
\voldef{2}{Geometria e Medidas}{Geometria e Medidas}
\capdef{v2c01}{2}{1}{Grandezas e unidades de medida}{Unidades de comprimento, área, volume, capacidade, massa, tempo; conversões; prefixos; avaliação de medições.}
\capdef{v2c02}{2}{2}{Escalas}{Escala em mapas, plantas e maquetes; razão entre áreas e volumes em escala.}
\capdef{v2c03}{2}{3}{Ângulos e posições de retas}{Ângulos, retas paralelas e perpendiculares, transversais, ângulos em polígonos, ângulos no relógio, giros e rotações em graus.}
\capdef{v2c04}{2}{4}{Figuras planas: polígonos e quadriláteros}{Características de figuras planas, triângulos, quadriláteros, polígonos regulares, ladrilhamentos, classificação.}
\capdef{v2c05}{2}{5}{Simetrias, transformações e localização}{Simetria axial e central, reflexão, rotação, translação, homotetia, localização e movimentação no plano e no espaço, trajetórias, coordenadas em mapas e malhas.}
\capdef{v2c06}{2}{6}{Congruência, semelhança e Teorema de Tales}{Congruência e semelhança de triângulos, razão de semelhança, Teorema de Tales, sombras e alturas.}
\capdef{v2c07}{2}{7}{Relações métricas nos triângulos}{Teorema de Pitágoras, relações métricas no triângulo retângulo, desigualdade triangular, lei dos senos e cossenos.}
\capdef{v2c08}{2}{8}{Trigonometria do ângulo agudo}{Seno, cosseno e tangente no triângulo retângulo, ângulos notáveis, rampas, inclinações, alturas inacessíveis.}
\capdef{v2c09}{2}{9}{Circunferência e círculo}{Elementos da circunferência, comprimento, arcos, setor circular, polígonos inscritos e circunscritos, rodas e engrenagens.}
\capdef{v2c10}{2}{10}{Perímetros e áreas}{Perímetro e área de figuras planas, figuras compostas, razão entre áreas, pisos e revestimentos, terrenos.}
\capdef{v2c11}{2}{11}{Sólidos geométricos, planificações e vistas}{Poliedros, relação de Euler, corpos redondos, planificações, vistas ortogonais, projeções, seções planas.}
\capdef{v2c12}{2}{12}{Volumes e áreas de superfície}{Volume e capacidade de prismas, cilindros, pirâmides, cones, esferas e sólidos compostos; área de superfície; embalagens; reservatórios.}
\voldef{3}{Álgebra, Funções e Geometria Analítica}{Álgebra, Funções e Geometria Analítica}
\capdef{v3c01}{3}{1}{Equações e inequações}{Equações e inequações do 1.º e do 2.º grau, modelagem algébrica de situações-problema, expressões algébricas.}
\capdef{v3c02}{3}{2}{Sistemas de equações}{Sistemas lineares 2x2 e 3x3, modelagem com duas ou mais incógnitas, interpretação de soluções.}
\capdef{v3c03}{3}{3}{Gráficos e funções}{Conceito de função, domínio e imagem, leitura e interpretação de gráficos cartesianos, crescimento e decrescimento, taxa de variação, funções definidas por partes.}
\capdef{v3c04}{3}{4}{Função afim}{Função polinomial do 1.º grau, coeficiente angular e linear, gráfico, modelagem de custos, tarifas e planos.}
\capdef{v3c05}{3}{5}{Função quadrática}{Função polinomial do 2.º grau, raízes, vértice, máximo e mínimo, parábolas, trajetórias.}
\capdef{v3c06}{3}{6}{Funções polinomiais e racionais}{Funções polinomiais de grau maior que 2, funções racionais, proporcionalidade inversa e hipérbole, assíntotas intuitivas.}
\capdef{v3c07}{3}{7}{Função exponencial}{Potências, função exponencial, crescimento e decaimento, juros como função, meia-vida, populações.}
\capdef{v3c08}{3}{8}{Logaritmos e função logarítmica}{Definição e propriedades de logaritmos, equações exponenciais via log, escalas logarítmicas (Richter, pH, decibéis).}
\capdef{v3c09}{3}{9}{Ciclo trigonométrico e funções trigonométricas}{Radianos, ciclo trigonométrico, seno e cosseno no ciclo, funções trigonométricas, período, amplitude, fenômenos periódicos.}
\capdef{v3c10}{3}{10}{Plano cartesiano e retas}{Coordenadas, distância entre pontos, ponto médio, equação da reta, coeficiente angular, paralelismo e perpendicularidade, regiões do plano.}
\capdef{v3c11}{3}{11}{Circunferência na geometria analítica}{Equação da circunferência, centro e raio, posições relativas entre ponto, reta e circunferência.}
